\section{Kolloquium}

Als Abschluss der Abschlussarbeit findet ein Seminarvortrag statt. Der Termin wird rechtzeitig bekannt gege-ben. 

Beim Seminarvortrag sind in der Regel dir Gutachter, Mitarbeiter, Studiernde  und andere interessierte Gäste zugegen. Du referieren zunächst in einem möglichst frei zu haltenden Kurzvortrag von nicht mehr als 15 Minuten Dauer über wesentliche Ergebnisse deiner Arbeit. Es folgt Diskussion und Fragerunde von ma-ximal 30- 45 Minuten. Im Anschluss an den Seminarvortrag ist nach Möglichkeit eine Vorführung der konkreten Ergebnisse erwünscht.

\subsection{Empfehlungen zum Kolloquiumsvortrag}

Für den Vortrag ist eine erneute Auswahl und Bewertung des in der schriftlichen Fassung der Diplomarbeit niedergelegten Stoffes notwendig. Durch den Kurzvortrag sollen Sie bei einem Zuhörerkreis, der im allge-meinen mit der Problematik Ihrer Arbeit nicht vertraut ist, Verständnis und Interesse für die Arbeit wecken und erzielte Ergebnisse erläutern. Eine didaktisch klare und straffe Gliederung Ihrer Überlegungen ist uner-lässlich. Ein Probevortrag vor Ihrem Betreuer wird dringend nahegelegt.

\susubsection{Folienvorlagen}

Eine PowerPoint- oder Tex-Vorlage für die Präsentation erhalten Sie auf Anfrage von Ihrem Betreuer. Zur Abschätzung der Anzahl der benötigten Folien gehen Sie bitte davon aus, dass pro Folie etwa 2 Min. Vor-tragszeit anzusetzen sind. 


Wenn die Zeitvorgaben für Vorträge um mehr als 2 min überzogen werden, wird der Vortrag abgebrochen. 

Die Bachelor- bzw. Masterarbeit wird durch das Abschlusskolloquium beendet. Es beginnt mit einer kurzenBegrüßung durch den Vorsitzenden der Prüfungskommission (i.Allg. also dem Erstgutachter) und
wird mit der Präsentation fortgesetzt, in der Sie die Ergebnisse Ihrer Arbeit vorstellen. Sofern Sie etwas vorführen möchten, dürfen Sie dies unmittelbar nach Ihrer Präsentation tun (dafür stehen ca. 5 min zur Verfügung). Es folgt die eigentliche Verteidigung, in der die Gutachter Sie zu Ihrer Arbeit befragen. Zum Abschluss berät sich die Prüfungskommission über die zu vergebende Note – natürlich unter Ausschluss der Öffentlichkeit. Sie erfahren das Ergebnis der Notenfindung unmittelbar nach der Beratung. Sofern Sie bestanden haben, bekommen Sie sofort eine Bescheinigung über den erfolgreichen Abschluss Ihres Studiums. Präsentation. Die Präsentation sollte für einen Beamer gestaltet und ca. 20 min lang sein. Sofern das Abschlusskolloquium in der HTW stattfindet, steht i.Allg. ein Whiteboard zur Verfügung, das Sie zur Beantwortung von Fragen nutzen können. Bereiten Sie sich auf die Präsentation gut vor. Machen Sie sich klar, dass eine Präsentation immer eine Art Verkaufsgespräch ist, in dem Sie sich gut präsentieren sollten.
Deshalb sollten Sie sich Gedanken zu folgenden Fragen machen:
 \begin{itemize}
     \item Wie sind Sie gekleidet?
\item Wie stehen Sie vor Ihrem Publikum?
\item Wie treten Sie auf?
\item Wie drücken Sie sich aus?
\item Sprechen Sie frei?
\item Wie sind Ihre Folien gestaltet (siehe unten)?
 \end{itemize}

Versuchen Sie, dem Vortrag eine klare Struktur zu geben. Er beginnt normalerweise mit einer kurzen Inhaltsübersicht, auf die eine knappe Einleitung folgt, in der Sie grundsätzlich in die Thematik Ihrer Abschlussarbeit einführen. Im Hauptteil der Präsentation beschreiben Sie die Resultate Ihrer Arbeit, wobei Sie davon ausgehen dürfen, dass Ihr Publikum fachkundig und mit dem Thema vertraut ist. Trotzdem sollten Sie Begriffe, soweit diese nicht zum allgemeinen Sprachgebrauch eines fachvertrauten Spezialisten zählen, kurz erläutern – besser jedoch vermeiden. Zum Ende des Vortrags empfiehlt sich ein Ausblick auf mögliche Verbesserungen, Erweiterungen usw.
Eine Zusammenfassung der Präsentation ist unnötig, da die Zeit dafür zu knapp ist und der Vortrag auch nicht durch Wiederholungen langweilen sollte. Vermeiden Sie, Ihre Präsentation einfach wortlos ausklingen zu lassen. Dadurch entstehen quälende Momente der Stille. Schließlich kann Ihr Publikum nicht wissen, ob mit dieser Pause Ihr Vortrag enden soll oder Sie nur eine kurze Bedenkzeit benötigen. Deshalb überlegen Sie sich im Vorfeld, wie Sie den Schluss gestalten möchten (z.B. „Ich danke Ihnen für Ihre Aufmerksamkeit“).
Abschließend noch ein Tipp: Stellen Sie sicher, dass am Tage des Abschlusskolloquiums keine technischen Probleme auftreten. So sind nicht in allen Räumen der HTW Beamer mit einem HDMI Eingang vorhanden. Auch sollten Sie wissen, wie Ihr Laptop zu konfigurieren ist, um die Folien über den Beamer auszugeben und wie der Bildschirmschoner ausgeschaltet werden kann, damit er nicht mitten im Vortrag aktiviert wird. Falls Sie eine Vorführung planen, müssen Sie sich außerdem darum kümmern, dass dafür alles notwendige bereit steht und auch lauffähig ist. 

\subsection{Folien} Es gibt zur Gestaltung von Folien eine große Zahl einfacher Regeln, deren wesentliche Zielsetzung es ist, zu übersichtlichen und ästhetischen Folien zu gelangen. Eine Faustregel ist maximal vier Worte pro Zeile und sechs Zeilen pro Folien zu verwenden (Ausnahmen bestätigen die Regel). Bevorzugen Sie auf
den Folien Stichpunkte statt ausgeschriebene Sätze. Dadurch vermeiden Sie, dass Ihr Publikum Texte liest statt Ihnen zuzuhören. Außerdem bewahren Sie den Zuhörer vor Langeweile. Sollte er nämlich bereits alles gelesen haben, was von Ihnen vorgetragen werden soll, schaltet er möglicherweise ab.

Achten Sie darauf, nicht zu viele Folien zu verwenden. Haben Sie z.B. 20 Folien geplant, dann steht Ihnen bei der vorgegebenen Vortragslänge von 20 Minuten gerade einmal 1 Minute pro Folie zur Verfügung. Das ist nicht viel. Bei sechs Zeilen pro Folie haben Sie zu jedem Stichpunkt gerade einmal 10 Sekunden Zeit – also gerade genug zum Vorlesen der Zeilen aber nicht für ein paar erklärende Worte. Verwenden Sie einen gut lesbaren und großen Schrifttyp. Vermeiden Sie den Einsatz von mehr als zwei unterschiedlichen Fonts auf einer Folie. Sonst wirken die Folien schnell unprofessionell und unruhig. Sie sollten auch auf „imposante“ Effekte verzichten. Diese stehen in PowerPoint zwar zur Verfügung, aber nur, damit sie nicht eingesetzt werden. Bewahren Sie Ihre Zuhörer vor einer endlosen aneinander Aneinanderreihung von Stichpunkten. Wie bei der Dokumentation gilt, dass Bilder aussagekräftiger sind, als Texte.

Bei Bildern und auch bei Tabellen (die ich hier nicht losgelöst betrachten möchte) sollten Sie darauf achten, diese nicht zu überladen. Arbeiten Sie jeweils die Informationen heraus, die für den Vortrag relevant sind. Das geht sehr einfach mit selbst gestalteten Bildern. Aber auch bei Bildern aus fremden Quellen
lassen sich Sachverhalte z.B. mit Hilfe einer gemalten Lupe hervorheben. Vernachlässigen Sie nicht den Gesamteindruck, den Ihre Folien vermitteln. So sollten Bilder selbstverständlich den Rand einer Folie nicht überragen. Problematisch ist auch die Verwendung von Screenshots. Durch die geringe Auflösung der meisten Beamer wirken solche Bilder schnell „pixelig“. Bei der Gestaltung Ihrer Folien ist ein ästhetisches Empfinden von Vorteil. Wenn Sie es nicht selbst haben, holen Sie den Rat einer Person guten Geschmacks ein. Ganz grundsätzlich: Die Folien sollten nicht im Auge weh tun. Bezogen auf eingebettete Bilder bedeutet dies, dass gedämpfte Farben Ton in Ton Verwendung finden (also ähnliche oder abgestufte Farben). Noch ein Tipp: Bei Fotos fällt manchmal negativ auf, wenn diese bei unterschiedlichen Lichtverhältnissen entstanden sind. Wenn Sie hier nicht aufwändig nacharbeiten möchten, können Sie sich behelfen, indem Sie die Farbe aus dem Bild heraus nehmen. Solche Fotos können Ihre Folien sogar aufwerten, weil Sie einen künstlerischen Aspekt einbringen. Abschließend sei angemerkt: Bei allen Sachverhalten, die Sie nicht selbst erarbeitet haben, müssen Sie Quellen benennen. Das gilt auch für Ihre Folien. Sie können hier genauso verfahren wie in Ihrer Dokumentation (siehe oben). Allerdings empfehle ich für die Folien, die Quellen immer dort anzugeben, wo Sie zitiert werden – z.B. in Kleinschrift direkt unter einem Bild. Auftreten. Versuchen Sie den Vortrag frei zu halten. Sollten Sie einen freien Vortrag, auch mit Übung vor
Freunden, Bekannten oder Verwandten nicht hin bekommen, dann machen Sie sich Kärtchen auf denen Stichpunkte stehen. Vermeiden Sie das reine Ablesen. Dabei kleben die Augen an Ihren Karten und es kommt zu keinem Augenkontakt mit den Zuhörern. Der ist jedoch wichtig, damit sich das Publikum angesprochen fühlt. Falls Sie zu ängstlich sind, um direkten Augenkontakt aushalten zu können, blicken Sie
leicht über oder neben Ihre Zuhörer. Allerdings funktioniert auch das nur, wenn Sie mehr als drei Zuhörer
haben.

Eine gewisse Nervosität ist bei einem öffentlichen Vortrag normal. Sie sollten sich davon nicht beeindrucken lassen. Nehmen Sie sich selbst gegenüber eine gewisse Trotzhaltung ein und sagen Sie sich, dass Sie die Situation beherrschen. Autosuggestion kann Ihnen tatsächlich helfen. Unterstützen Sie die so gewonnene
künstliche Selbstsicherheit, indem Sie eine entsprechende Körperhaltung einnehmen. Sie sollten mit durchgedrückten leicht geöffneten Beinen und aufrechtem Kopf gerade vor Ihrem Publikum stehen. Die Hände gehören dabei nicht in die Tasche. Eine Fernbedienung für Ihren Laptop kann wie das Brett in der See wirken, an dem Sie sich festhalten können. Bewegen Sie sich ruhig und gestikulieren Sie nicht zu viel.

Reden Sie laut und akzentuiert. Das schafft ebenfalls Selbstsicherheit. Außerdem werden Sie von Ihre Zuhörern verstanden. Beachten Sie bei allen Vorträgen, mit welchem Publikum Sie es zu tun haben. Bei Abschlussarbeiten sind das gewöhnlich Fachleute, die einen Vortrag erwarten, der wissenschaftlichen Ansprüchen genügt. Humor kann darin vorsichtig eingesetzt werden. Vermeiden Sie aber unbedingt, Posse auf Posse zu präsentieren. Fragen. Treten während Ihres Vortrags Fragen auf, so müssen Sie abwägen, wie Sie verfahren wollen. Es gibt Fragen, die sofort beantwortet werden müssen, weil sie essentiell für das Verständnis Ihres Vortrags sind. In den allermeisten Fällen können Sie die Fragen aber auch an das Ende Ihres Vortrags verschieben.